\section{Sharp four-option risks across sampling budgets}
\label{app:budgets}

\paragraph{A binary certificate for finite populations.}
Let there be $M$ possible orders, and let $\mathcal B$ be a nonempty
family of allowed $k$-subsets, $0<k<M$. A design samples an equally
weighted block with order inclusion probabilities $k/M$. For
$b_B=\mathbf 1_B$ write $u_B=(b_B-(k/M)\mathbf 1)/\sqrt{k(1-k/M)}$,
$W=\mathrm{span}\{u_B:B\in\mathcal B\}$, $r=\dim W$ and $T=M/k-1$.
Suppose a feasible design $D_*$ satisfies the weighted finite tight
frame identity $\mathbb E_{A\sim D_*}u_Au_A^\top=P_W/r$, where
$P_W$ is the orthogonal projector \citep{benedetto2003}. Its normalized
worst-case risk is $T/r$ for real profiles and for differences of
binary correctness indicators.

For any nonconstant $d$, let $z$ be its centered vector divided by its
Euclidean norm. The squared error of a realized block divided by
the full-order variance is $T\langle u_B,z\rangle^2$. For $D_*$,
averaging over $B$ gives $(T/r)\|P_Wz\|^2$, which proves the upper
bound. For the lower bound, take an independent $A\sim D_*$ and
assign model A the correctness vector $\mathbf 1_A$ and model B the
zero vector. The normalized centered profile is $u_A$. For every
feasible $B$, the average normalized squared error over $A$ is
$T u_B^\top(P_W/r)u_B=T/r$. Averaging over any candidate design on
$B$ leaves this value unchanged. Some fixed binary profile must
therefore realize risk at least $T/r$ for that design. The argument
allows the competing design to use arbitrary weights on its feasible blocks.

\paragraph{Enumeration and exact matrices.}
Index the 24 \emph{display} orders lexicographically, using four labels
numbered from zero. For a position-balanced $k$-block, each label
appears first $k/4$ times. Enumerate all choices of $k/4$ of the six
rows starting with each label, take their union and test the remaining
position, precedence and directed-adjacency counts. At $k=8$ this
gives 255, 99 and 15 feasible blocks in the three nested classes. At
$k=12$ it gives 640, 180 and 22. For each template draw a uniform
option relabeling. For position balance also permute slots uniformly.
For precedence and adjacency keep the slots or reverse them with equal
probability. The eight-call attaining mixtures are explicit in
Table~\ref{tab:eight-templates}. The twelve-call templates are given
in the accompanying machine-readable certificate.

\begin{table}[h]
\centering
\caption{Eight-call samplers. Each tuple lists the lexicographic indices of the eight displayed rows
before random relabeling. A uniform pair of distinct Williams blocks is an equivalent full-adjacency sampler.}
\label{tab:eight-templates}
\begin{tabular}{llr}
\toprule
Balance & Template & Selection probability \\
\midrule
Position & $(0,1,8,9,16,17,18,19)$ & $6/7$ \\
         & $(0,3,7,11,12,16,20,23)$ & $1/7$ \\
Position and precedence & $(0,1,6,10,13,17,22,23)$ & $2/11$ \\
         & $(0,2,7,10,13,17,21,22)$ & $2/11$ \\
         & $(0,3,7,11,12,16,20,23)$ & $1/11$ \\
         & $(0,3,7,11,14,17,18,22)$ & $6/11$ \\
All three & $(0,1,8,10,13,17,19,23)$ & $4/5$ \\
         & $(0,2,7,10,13,16,21,23)$ & $1/5$ \\
\bottomrule
\end{tabular}
\end{table}

For each attaining law form the rational matrix
$C=(M/k^2)\mathbb E[b_B b_B^\top]-\mathbf1\mathbf1^\top/M$.
The verifier checks its marginals, its trace $T$, the identity
$C^2=\lambda C$ and $Cu_B=\lambda u_B$ for \emph{every feasible}
block in the corresponding class. The range of $C$ lies in $W$
because $C=T\mathbb E[u_Bu_B^\top]$. The final check makes its range
equal to $W$. Hence $C=\lambda P_W$ and
$r=T/\lambda$. In increasing order of constraints the resulting ranks
are $(14,11,5)$ at eight and $(14,11,6)$ at twelve. The binary
certificate proves optimality in both profile classes. Simple random
sampling has constant risk $(M-k)/(k(M-1))$ on centered profiles.
Symmetrizing any unrestricted candidate over population relabelings
gives this law, proving the unrestricted lower bound for correctness
as well.

The complement of a balanced $k$-block satisfies the corresponding
constraints at $24-k$. The two centered sample means differ by the
factor $-k/(24-k)$ since their weighted sum is zero. This bijection
preserves uniform marginals and scales every profile risk by
$k^2/(24-k)^2$. It transfers the four- and eight-call optima to
20 and 16 calls. At 24 the target is observed exactly.

\paragraph{How the full-adjacency sampler changes with budget.}
The six four-call Williams blocks partition the population. Their
15 pairwise unions exhaust the admissible eight-call blocks, with
probability $1/15$ per union in the certified sampler. Writing the six
block means as a population of size six, the finite population
correction for sampling two distinct means is $(6-2)/(2(6-1))=2/5$.
At twelve calls there are also the even and odd permutation halves.
An attaining sampler selects a uniform triple of Williams blocks
with probability $5/6$ and a uniform parity half with probability $1/6$.
The two added blocks supply the sixth residual direction.

At eight calls, position features together with the constant have
rank ten, and precedence raises that rank to thirteen. Exact rational
elimination gives these ranks. The complementary dimensions match
the frame ranks 14 and 11. For full balance, every centered feasible
block is constant on each of the six Williams groups, a centered space
of dimension five. Hence the two projections of
Section~\ref{sec:budget} are precisely onto the two feasible-block
spans. For item $i$, the squared-error expectation equals
$d_i^\top C d_i/24$. Summing across independent items gives
Eq.~\eqref{eq:residual-threshold}.

\section{Replaying fixed-budget risks from the 24-order profiles}
\label{app:eight-replay}

The target is the difference of mean accuracies averaged over all 24
orders. For each question we compress the paired full-order correctness
difference into its positive and negative order masks. Questions with
identical masks are counted within stratum. These anonymous histograms
suffice for point risks and the paired item bootstrap. The three
matched panels use Qwen2.5-3B-Instruct and SmolLM3-3B. The separate
6,144-question HellaSwag comparison uses Qwen2.5-1.5B-Instruct and
Granite-3.3-2B-Instruct. It shares 227 questions with the matched
HellaSwag panel. We report it separately throughout.

For each block the integer score is the number of its positive
orders minus its number of negative orders. At $k=4,8$ the item
loss is $((24/k)\,\mathrm{score}-\sum_\pi d_i(\pi))^2/576$.
The replay sums the loss against the exact rational support weights.
As a second calculation at eight calls, it projects every item
using integer orthogonalization of position and precedence features
and the six Williams-group means. Both calculations agree exactly on
all balanced item risks, and the four-call risks and full-order gaps
match the earlier releases.

\begin{table}[h]
\centering
\caption{Conditional MSE of the paired accuracy gap in units of $10^{-6}$.
The first three columns use 992 questions each. The final column uses the separate smaller-model HellaSwag cohort.}
\label{tab:eight-all}
\begin{tabular}{@{}llrrrr@{}}
\toprule
Calls & Design & HellaSwag 3B & MMLU 3B & LogiQA 3B & HellaSwag 1.5B/2B \\
\midrule
4 & Unrestricted distinct        & 79.836 & 65.745 & 63.033 & 8.490 \\
4 & Position minimax             & 22.125 & 21.377 & 21.629 & 4.713 \\
4 & Precedence minimax           & 21.592 & 20.181 & 20.091 & 4.640 \\
4 & Williams                     & 21.645 & 20.869 & 19.625 & 4.586 \\
\midrule
8 & Unrestricted distinct        & 31.934 & 26.298 & 25.213 & 3.396 \\
8 & Position minimax             & 8.493 & 8.170 & 8.297 & 1.810 \\
8 & Precedence minimax           & 8.018 & 7.611 & 7.720 & 1.730 \\
8 & Two distinct Williams blocks & 8.658 & 8.348 & 7.850 & 1.834 \\
8 & Two disjoint position blocks & 8.773 & 8.477 & 8.565 & 1.870 \\
8 & Two disjoint precedence blocks & 8.328 & 7.879 & 7.887 & 1.798 \\
8 & Rotations and reversals      & 7.837 & 8.437 & 9.082 & 1.702 \\
\bottomrule
\end{tabular}
\end{table}

The disjoint-pair baselines choose uniformly among disjoint pairs
from the corresponding four-call minimax sampler. At eight calls
these baselines have worst-case ratios $1/6$, $3/14$ and $2/5$
for position, precedence and full adjacency, respectively. For the
rotation baseline we randomly relabel the four forward rotations and
their four reversed orders \citep{khan2026} per question. This law
partitions the population into three equiprobable blocks. A binary
profile constant on each block attains worst-case ratio one. It gives
the lowest observed point MSE among the tested eight-call designs on
both HellaSwag cohorts (Table~\ref{tab:eight-all}).

We computed 50,000 paired bootstrap resamples within each HellaSwag
domain and MMLU subject, holding stratum counts fixed. LogiQA was
resampled as one stratum. The same draws apply to all designs and
both budgets within each panel. The reported two-sided percentile
intervals describe sensitivity to item composition and have
panel-specific coverage without multiplicity adjustment. They do not
alter the exact conditional MSE on the observed items.

\begin{table}[h]
\centering
\caption{Change in the adjacency gain from four to eight calls, in percentage points.
The last two columns give the percent MSE reduction from optimizing at eight calls
versus uniformly pairing disjoint four-call minimax blocks in the same class.}
\begin{tabular}{@{}lrrr@{}}
\toprule
Panel & Gain change [95\% interval] & Position & Precedence \\
\midrule
HellaSwag 3B & $-7.74\;[-9.26,-6.23]$ & 3.19 & 3.73 \\
MMLU 3B & $-6.27\;[-7.80,-4.77]$ & 3.61 & 3.40 \\
LogiQA 3B & $-4.01\;[-5.48,-2.55]$ & 3.13 & 2.12 \\
HellaSwag 1.5B/2B & $-7.22\;[-7.73,-6.71]$ & 3.20 & 3.80 \\
\bottomrule
\end{tabular}
\end{table}

On the separate HellaSwag cohort the retained share was $48.21\%$.
The adjacency gain fell from $+1.17\%$ at four calls to $-6.05\%$
at eight (eight-call 95\% interval $[-7.29,-4.82]\%$).
The inference records were fixed before the eight-call study. We
specified the comparisons before computing the eight-call empirical
risks, with access to the published four-call outcomes.
